\section{尖点形式}\srcpage{52}
\emergencystretch=3em

在 Arinkin、Beraldo、Campbell、Chen、Færgeman、Gaitsgory、Lin、Raskin 与 Rozenblyum（2024）以及 Campbell、Chen、Færgeman、Gaitsgory、Lin、Raskin 与 Rozenblyum（2024）两篇文章之后，几何 Langlands 等价在 Eisenstein 部分已经知道，尚待处理的是尖点部分。这里，尖点部分指常数项函子的核；但人们也知道，它同时也是由从 $\LocSys_{\Gd}$ 张量到不可约 $\Gd$-局部系统的开子空间 $\LocSys^{\irred}_{\Gd}$ 上而得到的部分。因此，我们考虑函子
\[
\begin{aligned}
\mathbb{L}^{\mathrm{cusp}}_G\colon{}& D(\Bun_G)\\
&{}\otimes_{D_{\qc}(\LocSys_{\Gd})}D_{\qc}(\LocSys^{\irred}_{\Gd})\\
&{}\to D_{\qc}(\LocSys^{\irred}_{\Gd}).
\end{aligned}
\]
我们只简要勾勒总体策略。

首先，证明 $\mathbb{L}^{\mathrm{cusp}}_G$ 既有左伴随 $(\mathbb{L}^{\mathrm{cusp}}_G)^L$，又有右伴随 $(\mathbb{L}^{\mathrm{cusp}}_G)^R$，且二者同构。更确切地说，由于 $\mathbb{L}^{\mathrm{cusp}}_G$ 的源与靶都是自对偶范畴，对偶函子 $(\mathbb{L}^{\mathrm{cusp}}_G)^{\vee}$ 可看作反方向的函子。随后证明：在相差同一个扭的意义下，这个函子既是 $\mathbb{L}^{\mathrm{cusp}}_G$ 的左伴随，也是它的右伴随。利用 $\mathbb{L}_G$ 与 Whittaker 系数的相容性，证明它是右伴随；利用与 Kac--Moody 局部化的相容性，证明它是左伴随。这正是 Arinkin、Beraldo、Chen、Færgeman、Gaitsgory、Lin、Raskin 与 Rozenblyum（2024）的内容。

最核心的输入是 Færgeman 与 Raskin（2025）的定理：$\mathbb{L}^{\mathrm{cusp}}_G$ 保守（用到「尖点形式自动缓增」）。他们的定理大致是这样证明的：把 Whittaker 系数等同于 $T^*\Bun_G$ 中正则幂零轨迹处的微茎（这一断言的精确版本参见 Nadler 与 Taylor（2025）），然后利用层的微局部性质，说明对缓增层该微茎的消失蕴涵完全消失。\footnote{对这一步骤，可以归约到 $D(\Bun_G)$ 中具有受限变差的部分，从而归约到具幂零奇异支撑的层。若正则幂零轨迹处的微茎消失，则奇异支撑受到更强的约束，而结果表明这与缓增性矛盾。}

这意味着，只需理解复合 $\mathbb{L}^{\mathrm{cusp}}_G\circ(\mathbb{L}^{\mathrm{cusp}}_G)^R$，并证明它同构于恒等函子。由于 $\mathbb{L}^{\mathrm{cusp}}_G$ 是 $D_{\qc}(\LocSys^{\irred}_{\Gd})$-线性的，且有左伴随，可知该函子由与 $\LocSys^{\irred}_{\Gd}$ 上某个完美复形 $A_{\Gd}$ 作张量给出。此外，由与 Kac--Moody 局部化、进而与 opers 的联系，$A_{\Gd}$ 由 oper 结构空间的同调给出。这已经证明 $A_{\Gd}$ 连通，且其零次同调群以一般 oper 结构空间的连通分支为自由生成元。另一方面，由于 $\mathbb{L}^{\mathrm{cusp}}_G$ 的左、右伴随相同，可知 $A_{\Gd}$ 自对偶。又因它连通，实际上它是 0 次的一个向量丛。通过「以一般 oper 结构空间的连通分支为自由生成元」这一解释，它实际上是带联络的向量丛，从而是 $\LocSys^{\irred}_{\Gd}$ 上的局部系统。\srcpage{53}

Beraldo、Kazhdan 与 Schlank（2022）的近期定理表明：对经典群，一般 oper 结构空间非空且连通，于是经典群的证明完成。

另一方面，至少在 $g\ge 2$、并排除 $g=2$ 且 $G$ 为 $A_1$ 型的情形时\footnote{其余情形的尖点部分可以容易地手工处理。}，可以计算 $\LocSys^{\irred}_{\Gd}$ 的基本群；它与 $G$ 的中心的连通分支群有关。若 $G$ 的中心连通，则看出 $A_{\Gd}$ 是 $\LocSys^{\irred}_{\Gd}$ 的结构层的若干拷贝的直和；注意后者（即 $\LocSys^{\irred}_{\Gd}$）实际上是连通的。\footnote{这一连通性是另一个关键要素，且与算术 Langlands 的情形大不相同：在几何 Langlands 纲领中，尖点形式及其对偶 $L$-参数同样构成一个「连续谱」。}此外，$\LocSys^{\irred}_{\Gd}\subset\LocSys_{\Gd}$ 的补集余维至少为 2，且全空间是局部完全交；于是关于深度（depth）的考量表明，只需计算 $\LocSys_{\Gd}$ 上某个 $H^0$ 的维数。但把伴随展开，这由 Whittaker 层的自同态给出。由于 Whittaker 层完全显式，其自同态极易计算，且为一维，正如所愿。若 $G$ 的中心不连通，本质上同样的论证可行，只需一些额外的扭。

\begin{remark}
这一最后论证的关键之处在于：在 de Rham 情形，局部系统的叠只有极少的整体截面，因此最终可以把证明「处处重数一」的问题，转化为证明「整体截面层面重数一」。只有后一陈述适合在自守侧做这一计算。
\end{remark}

\emergencystretch=0pt